\subsection{The bent biplane as the abstract Kummer configuration}
\label{subsec:kummer-biplane-k3}

The $[16,6,6]_2$ bent-function code admits an explicit identification with the
classical Kummer $16_6$ incidence configuration.  Its sixteen minimum supports are
\[
 \mathcal B_a=\{x\in\mathbb F_2^4:q_0(x+a)=1\},\qquad
 q_0=x_0x_1+x_2x_3.
\]
The matrix
\[
 M=\begin{pmatrix}
 1&1&0&1\\
 1&1&1&0\\
 1&0&1&1\\
 1&1&0&0
 \end{pmatrix}\in GL(4,2)
\]
sends the six-point support of $q_0$ to the standard cross in a $4\times4$
array: the three other points in the distinguished row together with the three other
points in the distinguished column.  Consequently $M\mathcal B_a$ is the row/column
cross centred at $Ma$ for every $a$, providing an explicit labelled incidence
isomorphism with the standard row/column Kummer configuration.

Thus the point-derived doily split has a Kummer reading: fixing one node leaves six
incident blocks and ten nonincident blocks.  Under the established shortening these
become respectively the six doily ovoids and the ten grid-complement supports.  The
full automorphism group remains
\[
 2^4\!:\operatorname{Sp}(4,2)\cong2^4\!:\!S_6,
 \qquad |\operatorname{Aut}|=11520,
\]
matching the classical Kummer group $W(D_6)/\{\pm1\}$.

The $32$-vertex point/block Levi graph is furthermore explicitly isomorphic to the
folded six-cube
\[
 Q_6/\{x\sim x+111111\},
\]
with spectrum
\[
 \boxed{6^1\oplus2^{15}\oplus(-2)^{15}\oplus(-6)^1}.
\]
Primary Kummer-surface theory realizes the $16$ points and $16$ blocks as nodes and
tropes of a Kummer quartic, whose minimal resolution is a K3 surface.  The present
result is therefore a genuine incidence-level K3 bridge.  It does not yet choose a
specific projective quartic in the binary labels or construct a chain map from its
K3 divisor/cochain lattice into the $45$-point transport-curvature precomplex; those
remain separate realization problems.
